%-----------------------------------------------------------------------
\section{Our way}
\label{sec:ourway}
%-----------------------------------------------------------------------

The method follows directly from the two facts established in
Section~\ref{sec:fold}: the geometric quadrature is right wherever the map is
non-degenerate, and the band it cannot represent carries a finite amount of
power (Corollary~\ref{cor:integrable}). So we keep the geometric quadrature
everywhere, detect the band, and inside the band redistribute the band's own
measured energy according to the fold profile. Nothing is added to the image
that was not already in it under the geometric model; only its distribution
changes.

\subsection{The idea in one paragraph}
\label{sec:ours-idea}

Sample the surface, refract, splat, and compute $|\detJ|$ by finite differences
of the map \eqref{eq:raymap}. Threshold $|\detJ|$ to find the band. Outside the
band, keep \eqref{eq:irradiance} untouched. Inside the band, replace the
divergent weight with the peak-normalised Airy profile evaluated at the signed
distance to the band, and rescale that band's values so the band integrates to
the flux it had before. Because the correction is confined to the band and
conserves the band's own power, it cannot change the image outside the feature
and cannot introduce energy at all.

\subsection{Fold detection}
\label{sec:ours-detect}

The detector matters as much as the correction, and a careless one will appear
to work while being unable to converge.

\begin{lemma}[Fixed-quantile detectors cannot converge]
\label{lem:quantiledetector}
Let a detector flag a fixed quantile $p \in (0, 1)$ of any scalar field (whether image
radiance or Jacobian determinant $|\detJ|$). Then the measure of the flagged set is
fixed at $p \cdot \operatorname{area}(\recv)$ for all resolutions, independent of the
physical wavelength $\lambda$ or local wavefront curvature. Under mesh refinement,
the flagged set cannot contract to the true caustic curve $\Hm(\fold)$ (which has
Lebesgue measure zero).
\end{lemma}

\begin{proof}
By definition of the order statistic, the measure of the set where the field is below
its $p$-th quantile is $p \cdot \operatorname{area}(\recv)$ identically.
As the sampling resolution $N \to \infty$, the measure of the flagged domain remains
constant at $p$, rather than shrinking to zero as required for an asymptotic neighbourhood
of the 1D fold curve.
\end{proof}

Consequently, a physically consistent detector must threshold a geometric quantity
against a physical scale parameter:
\begin{equation}
B \;=\; \big\{\,\mathbf{x} \in \recv : |\detJ(\mathbf{x})| \le \epsilon\,\big\},
\label{eq:physicaldetector}
\end{equation}
where the threshold $\epsilon$ relates directly to the physical fringe width $w \propto \lambda^{2/3}$.
Transverse to the fold curve $\fold$ in the surface domain, let $x_{\perp}$ denote the
normal distance. By Taylor expansion around the critical set, the Jacobian determinant
scales as $|\detJ(\mathbf{x})| \approx |\nabla_{\perp}(\detJ)|\,|x_{\perp}|$.
On the receiver plane, the transverse distance from the caustic fold edge scales quadratically,
$b \approx \frac{1}{2} C_{\text{map}} x_{\perp}^2$. Because the physical Airy diffraction corridor
extends across receiver distance $b \lesssim C_w w \propto \lambda^{2/3}$, the corresponding
transverse boundary in the surface domain is $|x_{\perp}| \lesssim \sqrt{2 C_w w / C_{\text{map}}} \propto \lambda^{1/3}$.
Hence, the Jacobian threshold isolating the diffraction band is
$\epsilon = |\nabla_{\perp}(\detJ)|\,\sqrt{2 C_w w / C_{\text{map}}} \propto \lambda^{1/3}$.
Crucially, for a given scene and physical wavelength $\lambda$, $\epsilon$ is a \emph{fixed physical threshold}.
Holding $\epsilon$ fixed as the computational grid resolution $N$ increases ($N = 256 \to 2048$):
the physical area of the detected band stabilizes at $\operatorname{area}(B) \approx \SI{1.87}{\percent}$ of the receiver,
while the number of discrete grid cells inside the band grows quadratically ($1{,}311 \to 78{,}495$ cells),
accurately sampling the interior diffraction profile.
Conversely, in the geometric limit $\lambda \to 0$ (or $w \to 0$), the physical threshold $\epsilon \to 0$,
contracting the physical band area $\operatorname{area}(B) \to 0$ onto the measure-zero fold curve $\Hm(\fold)$.

In our discrete solver on the reference benchmark scene, we calibrate $\epsilon_0 = 0.0592$
from the 2nd percentile of $|\detJ|$ on the base grid ($p = 0.02$), which
reliably isolates the singular fold corridor across the receiver without referencing
image brightness. However, as Lemma~\ref{lem:quantiledetector} establishes, generalisation
to arbitrary scenes and scaling limits requires evaluating the absolute scale threshold \eqref{eq:physicaldetector}.

Section~\ref{sec:results-ablation} measures the alternative directly: swapping
the $|\detJ|$ detector for an image-brightness percentile detector degrades the
convergence of the result, because bright non-caustic glints are falsely flagged
while dimmer caustic fringes are missed.

\subsection{Distance transform and domain mapping}
\label{sec:ours-distance}

Once the fold band $B$ is flagged, we compute the distance $s$ using an exact
Euclidean distance transform (EDT) \cite{felzenszwalb2012}. 
Here, $s_i \ge 0$ is measured inwards from the non-band boundary toward the caustic core.
Because geometric transport carries essentially zero flux into the shadow zone ($G_i \approx 0$),
the wave correction inside $B$ maps $s_i$ to the illuminated oscillatory side of the
canonical fold profile, $u = s_i / w_c \ge 0$. The evaluation table spans $u \in [0, 8]$,
encompassing the primary caustic peak at $u = 1.0188$ and the first several diffraction fringes.
Outside $B$, the uncorrected geometric splat field is preserved unmodified.

\subsection{Normalisation, and why it is a discrete ratio}
\label{sec:ours-normalisation}

We now come to the point Lemma~\ref{lem:notl1} forces. The band's correction
multiplies a profile by a scale, and the scale is not a constant that can be
written down: because $|\Ai|^2$ is not integrable, there is no $Z$ with
$Z^{-1}|\Ai|^2$ a probability density, so a closed-form normaliser does not
exist. What does exist is a discrete ratio over the detected band, and
Corollary~\ref{cor:integrable} is what makes it meaningful.

Let $G_i$ be the geometric field \eqref{eq:irradiance} on the band $B$, let
$\bar G$ be its spatially smoothed version, let $\kappa$ be the ceiling,
let $b(s_i)$ be the profile \eqref{eq:profile}, and define
\begin{equation}
k^{*} \;=\;
\frac{\displaystyle \sum_{i \in B} G_i}
     {\displaystyle \sum_{i \in B} \min\!\big(\bar G_i, \kappa\big)\, b(s_i)} .
\label{eq:kstar}
\end{equation}

\begin{lemma}[Band-local flux is conserved]
\label{lem:flux}
With $I_i = \min(\bar G_i,\kappa)\,b(s_i)\,k^{*}$ for $i \in B$ and no ceiling
applied, the corrected band carries exactly the geometric band's flux:
$\sum_{i \in B} I_i = \sum_{i \in B} G_i$.
\end{lemma}

\begin{proof}
Substituting \eqref{eq:kstar}, and writing each step on its own line because the
chain does not fit a single column measure,
\begin{equation}
\begin{split}
\sum_{i \in B} I_i
&= k^{*} \sum_{i \in B} \min(\bar G_i,\kappa)\,b(s_i) \\[2pt]
&= \frac{\sum_{i\in B} G_i}{\sum_{i\in B}\min(\bar G_i,\kappa)b(s_i)}
   \cdot \sum_{i \in B} \min(\bar G_i,\kappa)\,b(s_i) \\[2pt]
&= \sum_{i \in B} G_i . \qedhere
\end{split}
\end{equation}
\end{proof}

The ceiling in the actual construction is
$I_i = \min\big(\min(\bar G_i,\kappa)b(s_i)k^{*},\,1.35\kappa\big)$, which can
clip the brightest band pixels and therefore breaks the exact equality of
Lemma~\ref{lem:flux}. We state the lemma for the unclipped case and report the
size of the discrepancy from a separate audit
(Section~\ref{sec:verify-energy}) rather than claiming an exactness the
implementation does not have.

The tempting alternative is to rescale the whole channel after the display clamp
so the overall level looks right. That fails: the clamp has already bounded the
image, and a subsequent global multiply pushes the peak back above the ceiling
the clamp just enforced, reintroducing the artefact the ceiling exists to
prevent. Conserving flux inside the band and applying the ceiling afterwards
keeps the bound at every resolution. This ordering is not cosmetic; it is the
difference between a bounded field and an unbounded one.

\begin{remark}[The ratio is not evidence]
\label{rem:circular}
Equation~\eqref{eq:kstar} divides the measured band flux by the measured band
flux by construction, so quoting $k^{*}$ as evidence of energy conservation is
circular. It is a bookkeeping step, not a measurement. The non-circular
statement is the throughput audit of Section~\ref{sec:verify-energy}, which
follows the light through every interface with the Fresnel weights of
\eqref{eq:fresnel}.
\end{remark}

\subsection{The algorithm}
\label{sec:ours-algorithm}

Algorithm~\ref{alg:bake} states the bake. The pass structure optimizes cost:
the surface evaluation, Snell refraction, and splatting are fused into one pass,
evaluating surface cells in a single cache-local sweep.

\begin{algorithm}[t]
\caption{The wave-corrected bake}
\label{alg:bake}
\begin{algorithmic}[1]
\Require surface $h$ \eqref{eq:surface}, source $\dir$, media $n(\lambda)$
         \eqref{eq:sellmeier}--\eqref{eq:rubin}, grid $N$, width $w_{\text{ref}}$,
         ceiling $\kappa$, threshold $\epsilon$
\Ensure field $I$ on the receiver
\Statex
\Statex \textbf{Pass 1 --- geometric backbone (fused)}
\ForAll{surface cells $\mathbf{x}$ on the grid}
    \State $\nrm \gets$ normal \eqref{eq:normal}; \quad $c_i \gets -\dir\cdot\nrm$
    \If{$\eta^2(1-c_i^2) > 1$} \textbf{continue} \Comment{total internal reflection}
    \EndIf
    \State $\trn \gets$ refracted direction \eqref{eq:snell}
    \State $\mathbf{u} \gets \Hm(\mathbf{x})$ \eqref{eq:raymap}
    \If{$\mathbf{u} \notin \recv$} \textbf{continue} \Comment{discard, do not clamp
        (Section~\ref{sec:ours-boundary})}
    \EndIf
    \State accumulate ray flux $\Delta\Phi$ (bilinear splat) into $G$ at $\mathbf{u}$
\EndFor
\Statex
\Statex \textbf{Pass 2 --- Jacobian and band}
\State $\detJ \gets$ finite differences of $\Hm$ \Comment{\eqref{eq:raymap}}
\State $B \gets \{\,i : |{\detJ}_i| \le \epsilon\,\}$ \Comment{Physical threshold \eqref{eq:physicaldetector}, calibrated to $\epsilon_0 = 0.0592$}
\State $s \gets$ Euclidean distance transform of $B$
\Statex
\Statex \textbf{Pass 3 --- per-channel Airy replacement}
\ForAll{channels $c \in \{R,G,B\}$}
    \State $w_c \gets w_{\text{ref}}\,(\lambda_c/550\,\text{nm})^{2/3}$
          \Comment{Corollary~\ref{cor:scaling}}
    \State $b(s) \gets 0.15 + 0.85\,|\Ai(-s/w_c)|^2 / \max_u|\Ai(-u)|^2$
          \Comment{\eqref{eq:profile}}
    \State $\bar G \gets$ spatially smoothed $G$; \quad
           $k^{*} \gets$ ratio \eqref{eq:kstar}
    \ForAll{$i \in B$}
        \State $I_i \gets \min\big(\min(\bar G_i,\kappa)\,b(s_i)\,k^{*},\;1.35\kappa\big)$
    \EndFor
\EndFor
\State \Return $I$
\end{algorithmic}
\end{algorithm}

\subsection{The receiver boundary}
\label{sec:ours-boundary}

Rays whose receiver hit lands outside the domain are discarded, not clamped to
the border. Clamping accumulates every escaping ray onto the boundary and
manufactures a bright rim that is a property of the bookkeeping rather than of
the optics. Formally, with $\recv$ the receiver footprint,
\begin{equation}
\text{contribute}(\mathbf{x}) \iff \Hm(\mathbf{x}) \in \recv .
\label{eq:boundary}
\end{equation}
The measured difference between the two choices is a factor of \num{38} in the
peak of an earlier version, and \eqref{eq:boundary} is also what makes the C++
and Python implementations agree to the precision reported in
Section~\ref{sec:verify-parity}.

\subsection{Coupling to a renderer}
\label{sec:ours-coupling}

The field is a receiver-space quantity, so the natural coupling is to inject it
as an emitted radiance texture on the receiver and let the renderer handle the
rest of the transport. Two details of that coupling cost us time and are worth
recording. The injection must be written as RGBA: a single-channel collapse
through an interpolation path silently discards the chromatic separation, and
the result looks plausible while being achromatic. And the injected texture must
not be used to compute the renderer's own caustics, or the two are counted
twice; the renderer's caustic pass is disabled for our frames.

\subsection{Cost and complexity}
\label{sec:ours-cost}

The bake is $\order{N^2}$ in the grid for pass 1, plus the cost of the distance
transform in pass 2, plus three channel passes in pass 3. The measured figures
are in Table~\ref{tab:cost}; the salient ones are that a $512^2$ bake takes
\SI{36.2}{\milli\second} on eight threads and that two thirds of it is the Airy
term, which is the obvious target for further work.

We deliberately do not claim the method is ``faster than path tracing'' in
general. The cost comparison that is meaningful is per caustic of a given
quality, and Section~\ref{sec:results-cost} makes it on that basis, including
the regime where we lose.

\subsection{The parameters}
\label{sec:ours-parameters}

The method has three parameters, and Section~\ref{sec:results-ablation}
measures the sensitivity to each rather than asserting a value.

\begin{itemize}
\item \textbf{Detector percentile.} Second percentile of $|\detJ|$. This sets
      the band width; too small and the band fails to cover the fold, too large
      and the correction is applied where the geometry was correct.
\item \textbf{Reference width $w_{\text{ref}}$.} The fringe scale
      \eqref{eq:scalelaw} for the geometry. Its physically meaningful value for
      this scene is \SI{123}{\micro\meter}
      (Sections~\ref{sec:estimators} and~\ref{sec:dispersion}), and using a
      larger value for visibility is exactly the exaggeration this paper avoids.
\item \textbf{Ceiling $\kappa$.} Bounds the field. It interacts with
      Lemma~\ref{lem:flux} as described, and the audit of
      Section~\ref{sec:verify-energy} quantifies the interaction.
\end{itemize}
